% !TEX root = userguide.tex

\def\headdate{31st January 2012}

\section{Grid deformation}
\label{sec:grid_deformation}

Starting from a base mesh elements can be moved/changed in size according to
a monitor function $f$. The value of $f$ gives the relative size of elements
at that position relative to the size of the other elements. Hence the absolute
value of $f$ is not important. A possibility for generating
the function $f$ is using a levelset function.
The method is described in \cite{Wan07,Grajewski09} and requires solving a
Poisson problem and a pseudo time stepping procedure. Is is used extensively in
the PhD thesis of Choi \cite{Choi11b}.

The example files are \texttt{grid\_deformation1.f90} \dots
\texttt{grid\_deformation6.f90},
and can be obtained by typing at the
command line:
\medskip
\begin{quote}
   \texttt{getexample grid\_deformation}
\end{quote}
The examples \texttt{grid\_deformation1.f90}  \dots
\texttt{grid\_deformation5.f90} describe the use of grid deformation for a
circle and sphere with fixed domain boundaries,
circle and sphere with periodical
boundaries in one direction and two circles with fixed domain boundaries. The
result of the latter example is shown in Figure~\ref{fig:griddefcircles}.
\begin{figure}[htbp!]
   \begin{center}
   \includegraphics[scale=0.4]{mesh_initial}
   \includegraphics[scale=0.4]{mesh_deformed}
   \end{center}
   \caption{Mesh before and after grid deformation.}
    \label{fig:griddefcircles}
\end{figure}

The example \texttt{grid\_deformation6.f90} is similar to the
\texttt{stokes8.f90} example (see Sec.~\ref{sec:drivenparticle}),
 but now with grid deformation.
